\documentclass[11pt]{article}

\usepackage[T1]{fontenc}
\usepackage[utf8]{inputenc}
\usepackage{lmodern}
\usepackage[margin=1in]{geometry}
\usepackage{amsmath}
\usepackage{amssymb}
\usepackage{amsthm}
\usepackage{array}
\usepackage{bbm}
\usepackage{booktabs}
\usepackage[labelfont=bf]{caption}
\usepackage{float}
\usepackage{graphicx}
\usepackage{hyperref}
\usepackage{longtable}
\usepackage{makecell}
\usepackage{multirow}
\usepackage[round,authoryear]{natbib}
\usepackage{url}

\graphicspath{{build/figures/}}
\hypersetup{
  colorlinks = true,
  linkcolor = blue,
  citecolor = blue,
  urlcolor = blue
}

\setlength{\parindent}{0pt}
\setlength{\parskip}{0.75em}
\setlength{\emergencystretch}{2em}
\hfuzz=3pt
\hbadness=10000
\allowdisplaybreaks

\DeclareMathOperator*{\E}{E}
\DeclareMathOperator*{\logit}{logit}
\DeclareMathOperator*{\Bernoulli}{Bernoulli}
\DeclareMathOperator*{\Beta}{Beta}

\newtheorem{proposition}{Proposition}
\newtheorem{lemma}[proposition]{Lemma}

\input{build/generated-values.tex}

\title{Additional analyses for two-component treatment-effect testing}
\author{
  Andreas Kryger Jensen and Theis Lange\\
  Section of Biostatistics, Department of Public Health\\
  University of Copenhagen
}
\date{}

\begin{document}

\maketitle

\section*{Introduction}

This document contains two additional analyses illustrating the two-component testing framework for outcomes with a clinically meaningful atom and a non-atom continuous or positive-score component. The examples are secondary to the IST-3 application in the main manuscript: the joineR liver data provide a public landmark-analysis example with death and a positive biomarker component, and the licorice gargle randomized trial provides a public atom-plus-positive-score example. They are included here to demonstrate the workflow on reproducible package data sets rather than to serve as primary clinical applications.

\section{joineR liver data}\label{app:joiner-liver-application}

As a secondary example, we applied the same workflow to the liver cirrhosis trial data distributed with the \texttt{joineR} R package, version \LiverAppendixPackageVersion{} \citep{joineRPackage}. The data contain repeated prothrombin index measurements, survival time, censoring status, and the allocated treatment arm for \LiverAppendixRandomizedN{} participants from a controlled trial of prednisone in liver cirrhosis. The package documentation states that \texttt{treatment = 0} denotes placebo, \texttt{treatment = 1} denotes prednisone, and \texttt{cens = 1} denotes death; the data were analyzed by \citet{henderson2002identification} and are described as coming from \citet{andersen2003statistical}. We use this example as an additional analysis because it is public and reproducible through an R package, but it requires a landmark construction and is therefore less direct than the IST-3 quality-of-life endpoint.

The secondary endpoint uses a two-year landmark. Participants who died before or at two years are assigned the atom value 0. Participants known alive at two years are assigned their latest recorded prothrombin index at or before two years. Participants censored before two years are excluded because their two-year death status is not observed. Thus
\[
  Y =
  \begin{cases}
    0, & \text{death by 2 years},\\
    \text{latest prothrombin index before 2 years}, & \text{known alive at 2 years},
  \end{cases}
  \qquad
  A = \mathbbm{1}(Y \ne 0).
\]
We code $R=1$ for prednisone and $R=0$ for placebo, so all contrasts are prednisone minus placebo. Higher prothrombin index is treated as the favorable direction for the continuous component, and avoiding the death atom is favorable for the binary component. This left \LiverAppendixN{} participants for analysis after excluding \LiverAppendixExcludedCensored{} participants censored before the landmark (\LiverAppendixExcludedControl{} placebo and \LiverAppendixExcludedTreatment{} prednisone). No additional subject was excluded for lacking a prothrombin measurement before the landmark.

\input{build/tables/appendix-liver-descriptives.tex}

Figure~\ref{fig:appendix-liver} displays the endpoint components and the semiparametric likelihood-ratio region. The death atom was common in both arms: \LiverAppendixControlDeaths{} placebo participants (\LiverAppendixControlDeathPct{}) and \LiverAppendixTreatmentDeaths{} prednisone participants (\LiverAppendixTreatmentDeathPct{}) died by two years among the landmark analysis set. Among participants known alive at two years, the mean prothrombin index was \LiverAppendixControlMean{} in the placebo arm and \LiverAppendixTreatmentMean{} in the prednisone arm, giving $\hat\mu_\delta=\LiverAppendixMuDelta{}$. The fitted odds ratio for being in the non-atom component rather than at the death atom was $\hat\alpha_\delta=\LiverAppendixAlphaDelta{}$, and the combined mean contrast under the death-at-zero coding was $\hat\Delta=\LiverAppendixDelta{}$.

\begin{figure}[htbp]
\centering
\includegraphics[width=\textwidth]{appendix-liver.pdf}
\caption{joineR liver endpoint components and semiparametric likelihood-ratio region. The left panel shows death by two years and known-alive-with-prothrombin counts by treatment arm. The middle panel shows the prothrombin distribution among known two-year survivors. The right panel shows the semiparametric likelihood-ratio surface for $(\mu_\delta,\log\alpha_\delta)$ with reference lines at no prothrombin mean difference and no non-atom probability difference.}
\label{fig:appendix-liver}
\end{figure}

Table~\ref{tab:liver-frequentist} compares standard analyses with the proposed likelihood-ratio tests. A Welch $t$-test applied to the atom-coded endpoint gave p-value \LiverAppendixTTestP{}, and the Wilcoxon rank-sum test gave p-value \LiverAppendixWilcoxP{}. As in the main application, the atom-coded $t$-test depends on the numerical atom value, whereas the Wilcoxon test treats deaths as tied at the bottom of the analyzed scale. Analyses restricted to the non-atom component gave p-values \LiverAppendixNonAtomTTestP{} for Welch's $t$-test and \LiverAppendixNonAtomWilcoxP{} for the Wilcoxon test, but those analyses condition on being known alive at the two-year landmark.

\input{build/tables/appendix-liver-frequentist.tex}

The parametric LRT statistic was \LiverAppendixLRTStatistic{} with p-value \LiverAppendixLRTP{}, and the semiparametric LRT statistic was \LiverAppendixSPLRTStatistic{} with p-value \LiverAppendixSPLRTP{}. Both tests target the joint null of no prednisone-placebo difference in death by two years and no difference in the prothrombin mean among known two-year survivors. In this example, the joint result is driven primarily by the continuous prothrombin component, while the estimated non-atom odds ratio is close to one. The result should therefore be read as an illustration of the method on a public clinical data set, not as a definitive treatment conclusion.

\input{build/tables/appendix-liver-delta-ci.tex}

Table~\ref{tab:liver-delta-ci} gives the Welch, profile, and projected intervals for $\Delta$. The same interpretational caveat applies as in the IST-3 analysis: $\Delta$ is a mean of the numerical combined endpoint after assigning death the value 0, so changing the atom coding would change its scale. The component estimates $(\mu_\delta,\alpha_\delta)$ remain directly interpretable as a prothrombin mean difference and a non-atom odds ratio.

We also fit the Bayesian two-part model in \texttt{TruncComp2}. The Bernoulli atom probabilities used independent $\mathrm{Beta}(1,1)$ priors. Because all non-atom prothrombin values are strictly positive, the continuous component used the positive-real mixture model with four components and weakly informative priors on the internally transformed scale. The fit is cached with the generated data; setting \texttt{TRUNCCOMP\_REFRESH\_LIVER\_BAYES=true} regenerates it from the packaged \texttt{joineR::liver} data.

\input{build/tables/appendix-liver-bayes-summary.tex}

The Bayesian summaries in Table~\ref{tab:liver-bayes} again separate the atom and continuous components. Under the benefit-oriented coding used here, $\Pr(\Delta>0\mid\mathrm{data})=\LiverAppendixDeltaBenefitProbability{}$ summarizes the posterior probability that prednisone improves the combined two-year endpoint. The corresponding posterior probabilities for the component effects are $\Pr(\mu_\delta>0\mid\mathrm{data})=\LiverAppendixMuBenefitProbability{}$ for the prothrombin mean and $\Pr(\alpha_\delta>1\mid\mathrm{data})=\LiverAppendixAlphaBenefitProbability{}$ for avoiding the atom. Figure~\ref{fig:appendix-liver-posterior} shows these posterior contrasts.

\begin{figure}[htbp]
\centering
\includegraphics[width=0.92\textwidth]{appendix-liver-posterior.pdf}
\caption{Posterior distributions for the joineR liver contrasts. Vertical dashed lines mark no prothrombin mean difference, no difference in the odds of being known alive with prothrombin observed, and no combined endpoint mean difference.}
\label{fig:appendix-liver-posterior}
\end{figure}

Finally, Table~\ref{tab:liver-ppc} and Figure~\ref{fig:appendix-liver-ppc} report posterior predictive checks for the atom and continuous components. These checks are especially important here because the endpoint uses a derived landmark biomarker summary rather than a trial-scheduled single assessment. The atom check assesses whether the model reproduces the two-year death proportions, while the continuous check assesses whether the positive-real mixture captures the skewness and spread of prothrombin among known two-year survivors. We report the sampler and truncation diagnostics in Table~\ref{tab:liver-bayes}; any remaining warnings should be interpreted as limitations of the experimental Bayesian mixture fit rather than as part of the frequentist likelihood-ratio analysis.

\input{build/tables/appendix-liver-ppc.tex}

\begin{figure}[htbp]
\centering
\includegraphics[width=\textwidth]{appendix-liver-ppc.pdf}
\caption{Posterior predictive checks for the joineR liver Bayesian model, separately for the two-year death atom and the prothrombin component among known two-year survivors.}
\label{fig:appendix-liver-ppc}
\end{figure}

\clearpage

\section{Licorice gargle randomized trial}\label{app:licorice-gargle-application}

As a second secondary example, we analyzed the licorice gargle randomized trial data distributed with the \texttt{medicaldata} R package, version \LicoriceAppendixPackageVersion{} \citep{medicaldataPackage}. The data source is the Licorice Gargle Dataset contributed by \citet{licoriceGargleDataset}; the source clinical trial is the randomized, double-blind comparison of preoperative licorice versus sugar-water gargle for prevention of postoperative sore throat and postextubation coughing after double-lumen endotracheal intubation \citep{ruetzler2013licorice}. The package data are reproducible from CRAN, but the original TSHS portal is an educational repository and states that publication reuse permission is not granted by the portal itself. We therefore present this as a transparent secondary illustration rather than as the main clinical application.

The scientific atom is absence of postoperative sore-throat pain. We used the 30-minute PACU pain-on-swallowing score, \texttt{pacu30min\_swallowPain}, because the documentation describes it as an 11-point Likert score with 0 denoting no pain and 10 denoting worst pain. The endpoint is
\[
Y_i =
\begin{cases}
0, & \text{if participant } i \text{ reported no swallowing pain at 30 minutes},\\
\text{positive swallowing-pain score}, & \text{otherwise},
\end{cases}
\]
with atom value 0 and $A_i = 1(Y_i \ne 0)$. We coded $R=1$ for licorice gargle and $R=0$ for sugar-water gargle, so all contrasts are licorice minus sugar-water. In contrast to the IST-3 and liver examples, larger values of $Y$ are worse. A clinically favorable effect is therefore a higher no-pain atom probability, a lower odds of being in the positive-pain component, a negative $\mu_\delta$, and a negative combined mean contrast $\Delta$.

\input{build/tables/appendix-licorice-descriptives.tex}

The packaged data contain \LicoriceAppendixRandomizedN{} participants, with \LicoriceAppendixExcludedMissing{} missing 30-minute swallowing-pain scores (\LicoriceAppendixExcludedControl{} in the sugar-water arm and \LicoriceAppendixExcludedTreatment{} in the licorice arm). The analysis includes \LicoriceAppendixControlN{} sugar-water participants and \LicoriceAppendixTreatmentN{} licorice participants. The no-pain atom was observed for \LicoriceAppendixControlNoPain{} sugar-water participants (\LicoriceAppendixControlNoPainPct{}) and \LicoriceAppendixTreatmentNoPain{} licorice participants (\LicoriceAppendixTreatmentNoPainPct{}). Among participants with positive pain, the mean score was \LicoriceAppendixControlMean{} in the sugar-water arm and \LicoriceAppendixTreatmentMean{} in the licorice arm, giving $\hat\mu_\delta=\LicoriceAppendixMuDelta{}$. The fitted odds ratio for any pain rather than no pain was $\hat\alpha_\delta=\LicoriceAppendixAlphaDelta{}$, and the combined mean contrast was $\hat\Delta=\LicoriceAppendixDelta{}$. These estimates all favor licorice on the chosen scale.

\begin{figure}[htbp]
\centering
\includegraphics[width=\textwidth]{appendix-licorice.pdf}
\caption{Licorice gargle endpoint components and semiparametric likelihood-ratio region. The left panel shows no-pain and any-pain counts by randomized arm. The middle panel shows the positive 30-minute swallowing-pain scores among participants with symptoms. The right panel shows the semiparametric likelihood-ratio surface for $(\mu_\delta,\log\alpha_\delta)$ with reference lines at no positive-pain mean difference and no difference in the odds of any pain.}
\label{fig:appendix-licorice}
\end{figure}

Table~\ref{tab:licorice-frequentist} compares standard analyses with the likelihood-ratio tests. A Welch $t$-test on the atom-coded pain endpoint gave p-value \LicoriceAppendixTTestP{}, and the Wilcoxon rank-sum test gave p-value \LicoriceAppendixWilcoxP{}. These standard tests are easy to compute, but the $t$-test remains tied to the numerical atom coding and the Wilcoxon test treats no-pain observations as tied at the bottom of a scale on which larger values are worse. The positive-pain-only Welch and Wilcoxon analyses gave p-values \LicoriceAppendixNonAtomTTestP{} and \LicoriceAppendixNonAtomWilcoxP{}, respectively, but condition on the post-randomization subset with symptoms. Fisher's exact test for the no-pain atom gave p-value \LicoriceAppendixAtomFisherP{}.

\input{build/tables/appendix-licorice-frequentist.tex}

The parametric LRT statistic was \LicoriceAppendixLRTStatistic{} with p-value \LicoriceAppendixLRTP{}, and the semiparametric LRT statistic was \LicoriceAppendixSPLRTStatistic{} with p-value \LicoriceAppendixSPLRTP{}. Both tests reject the joint null of no licorice-sugar-water difference in the no-pain probability and no difference in positive swallowing-pain scores. In this example the component estimates point in the same clinical direction: licorice is associated with more participants at the no-pain atom and lower positive pain scores among those with symptoms.

\input{build/tables/appendix-licorice-delta-ci.tex}

Table~\ref{tab:licorice-delta-ci} reports the Welch, profile, and projected intervals for $\Delta$. Because $Y=0$ is a favorable atom and larger positive scores are worse, intervals below zero favor licorice for the coded combined endpoint. This interpretation is clinically natural here, but it is still a numerical composite: changing the atom coding or reversing the pain scale would change $\Delta$, whereas $(\mu_\delta,\alpha_\delta)$ retain their component interpretations.

We also fit the Bayesian two-part model in \texttt{TruncComp2}. The Bernoulli atom probabilities used independent $\mathrm{Beta}(1,1)$ priors. Since all non-atom scores are strictly positive, the continuous component used the positive-real mixture model with two components and weakly informative priors on the internally transformed scale. The fit is cached with the generated data and can be refreshed from the packaged \texttt{medicaldata} data using the manuscript asset script.

\input{build/tables/appendix-licorice-bayes-summary.tex}

The Bayesian summaries in Table~\ref{tab:licorice-bayes} give the same clinical picture. Because larger pain scores are worse, the posterior probability of treatment benefit for the combined endpoint is $\Pr(\Delta<0\mid\mathrm{data})=\LicoriceAppendixDeltaBenefitProbability{}$. The corresponding component probabilities are $\Pr(\mu_\delta<0\mid\mathrm{data})=\LicoriceAppendixMuBenefitProbability{}$ for lower positive pain scores, $\Pr(\alpha_\delta<1\mid\mathrm{data})=\LicoriceAppendixAlphaBenefitProbability{}$ for lower odds of any pain, and $\Pr(\text{no-pain probability increase}\mid\mathrm{data})=\LicoriceAppendixNoPainBenefitProbability{}$. Figure~\ref{fig:appendix-licorice-posterior} shows the posterior contrasts with reference values for no effect.

\begin{figure}[htbp]
\centering
\includegraphics[width=0.92\textwidth]{appendix-licorice-posterior.pdf}
\caption{Posterior distributions for the licorice gargle contrasts. Vertical dashed lines mark no positive-pain mean difference, no difference in the odds of any pain, and no combined endpoint mean difference. Negative values favor licorice for $\mu_\delta$ and $\Delta$, and negative log odds ratios favor licorice for $\alpha_\delta$.}
\label{fig:appendix-licorice-posterior}
\end{figure}

Finally, Table~\ref{tab:licorice-ppc} and Figure~\ref{fig:appendix-licorice-ppc} report posterior predictive checks. The atom check assesses whether the model reproduces the no-pain proportions in each arm. The continuous check is more demanding because the positive component is an integer-valued, bounded Likert score with visible heaping rather than a truly continuous measurement. We therefore regard this analysis as useful for demonstrating the atom-plus-positive-score workflow, while the posterior predictive diagnostics are a reminder that the smooth positive-real Bayesian mixture is only an approximation for this endpoint.

\input{build/tables/appendix-licorice-ppc.tex}

\begin{figure}[htbp]
\centering
\includegraphics[width=\textwidth]{appendix-licorice-ppc.pdf}
\caption{Posterior predictive checks for the licorice gargle Bayesian model, separately for the no-pain atom and the positive 30-minute swallowing-pain score among participants with symptoms.}
\label{fig:appendix-licorice-ppc}
\end{figure}

\clearpage
\bibliographystyle{plainnat}
\bibliography{bibliography}

\end{document}
